\subsection{GROUP OF TANGENT VECTORS TO A LIE GROUP}

PROPOSITION 4. Let G be a Lie group. Then T(G), with the law of composition tangent to the multiplication of G, is a Lie group. The identity element of T(G) is the vector \(0_{e}\) .

This follows from Propositions 1 and 2.

PROPOSITION 5. Let G and H be Lie groups and f a morphism of G into H. Then T(f) is a morphism of the Lie group T(G) into the Lie group T(H).

We know that \(\mathrm{T}(f)\) is analytic. On the other hand, let m (resp. n) denote the multiplication on G (resp. H). Then \(f \circ m = n \circ (f \times f)\) , whence

\[
\mathrm{T} (f) \circ \mathrm{T} (m) = \mathrm{T} (n) \circ (\mathrm{T} (f) \times \mathrm{T} (f)),
\]

which expresses the fact that \(T(f)\) is a group homomorphism.

COROLLARY. Let \(G_{1},\ldots,G_{n}\) be Lie groups. The canonical isomorphism of the manifold \(\mathrm{T}(\mathrm{G}_{1}\times\cdots\times\mathrm{G}_{n})\) onto the manifold \(\mathrm{T}(\mathrm{G}_{1})\times\cdots\times\mathrm{T}(\mathrm{G}_{n})\) is a Lie group isomorphism.

\(pr_{i}\) is a morphism of \(G_{1} \times \cdots \times G_{n}\) into \(G_{i}\) and hence \(T(pr_{i})\) is a morphism of \(T(G_{1} \times \cdots \times G_{n})\) into \(T(G_{i})\) .

PROPOSITION 6. Let G be a Lie group.

\begin{enumerate}
\def\labelenumi{(\roman{enumi})}
\item
  The canonical projection \(p: \mathrm{T}(\mathrm{G}) \to \mathrm{G}\) is a Lie group morphism.
\item
  The kernel of \(p\) is \(\mathrm{T}_e(\mathrm{G})\) . It is a Lie subgroup of \(\mathrm{T}(\mathrm{G})\) . The Lie group structure induced on \(\mathrm{T}_e(\mathrm{G})\) by that on \(\mathrm{T}(\mathrm{G})\) is the Lie group structure of the complete normable space \(\mathrm{T}_e(\mathrm{G})\) .
\item
  The zero section \(s\) is an isomorphism of the Lie group \(\mathbf{G}\) onto a Lie subgroup \(s(\mathbf{G})\) of \(\mathrm{T}(\mathbf{G})\) (which subgroup we identify with \(\mathbf{G}\) ).
\item
  The Lie group T(G) is the semi-direct product of G by \(T_{e}(G)\) .
\end{enumerate}

Assertion (i) follows from (5). Assertion (ii) is obvious taking account of Proposition 2 (ii). Assertions (iii) and (iv) follow from (6) and § 1, Proposition 8.

Let \(u \in \mathrm{T}(\mathrm{G})\) and \(g \in \mathrm{G}\) . By (3) and (4), the products \(ug\) , \(gu\) calculated in the group \(\mathrm{T}(\mathrm{G})\) are the images of \(u\) under \(\mathrm{T}(\delta(g^{-1}))\) and \(\mathrm{T}(\gamma(g))\) . It follows from § 1, Corollary 2 to Proposition 17 that the mapping \((g, u) \mapsto gu\) of \(\mathrm{G} \times \mathrm{T}_e(\mathrm{G})\) into \(\mathrm{T}(\mathrm{G})\) is an isomorphism of the trivial vector bundle \(\mathrm{G} \times \mathrm{T}_e(\mathrm{G})\) with base space \(G\) onto the vector bundle \(\mathrm{T}(\mathrm{G})\) . The inverse isomorphism is called the left trivialization of \(\mathrm{T}(\mathrm{G})\) . By considering the mapping \((g, u) \mapsto ug\) , the right trivialization of \(\mathrm{T}(\mathrm{G})\) is defined similarly.

PROPOSITION 7. Let G be a Lie group, M a manifold of class \(\mathbf{C}^r\) and \(f\) and \(g\) mappings of class \(\mathbf{C}^r\) of M into G, so that \(fg\) is a mapping of class \(\mathbf{C}^r\) of M into G. Let \(m \in \mathbf{M}\) , \(x = f(m)\) , \(y = g(m)\) , \(u \in \mathrm{T}_m(\mathbf{M})\) . Then

\[
(\mathrm{T} f g) u = \mathrm{T} (f) u. y + x. \mathrm{T} (g) u.
\]

Let \(m\) be the multiplication of G. Then \(fg = m \circ (f, g)\) . Now

\[
\mathrm{T} (f, g) (u) = (\mathrm{T} (f) u, \mathrm{T} (g) u),
\]

hence \(\mathrm{T}(fg)u = \mathrm{T}(f)u.\mathrm{T}(g)u\) . It then suffices to apply (1) with \(f\) replaced by \(m\) .

COROLLARY. Let \(n \in \mathbf{Z}\) . The tangent mapping at \(e\) to the mapping \(g \mapsto g^n\) of \(G\) into \(G\) is the mapping \(x \mapsto nx\) of \(T_e(G)\) into \(T_e(G)\) .

For \(n \geqslant 0\) , this follows by induction on \(n\) from Proposition 7. On the other hand, the tangent mapping at \(e\) to the mapping \(g \mapsto g^{-1}\) is the mapping \(x \mapsto -x\) (no. 1, Proposition 2).

Let G be a Lie group, X a manifold of class \(C^r\) and \((g, x) \mapsto gx\) a law of left operation of class \(C^r\) of G on X. Arguing as for Proposition 1, we derive a law of left operation of class \(C^{r-1}\) of T(G) on T(X), which we shall also denote by \((u, v) \mapsto uv\) . Identifying G (resp. X) with the image of the zero section of T(G) (resp. T(X)), we see by (6) that the law of left operation of T(G) on T(X) extends the law of left operation of G on X. For all \(u \in \mathrm{T}_{g}(G)\) and \(v \in \mathrm{T}_{x}(X)\) , by (1),

\[
u v = g v + u x.\tag{7}
\]

If \(g \in G\) and \(v \in T_x(X)\) , \(gv\) is, by (3), the image of \(v\) under the tangent mapping at \(x\) to the mapping \(y \mapsto gy\) of \(X\) into \(X\) . This tangent mapping is an isomorphism of \(T_x(X)\) onto \(T_{gx}(X)\) . In particular,

\[
g (v + v ^ {\prime}) = g v + g v ^ {\prime}, \quad g (\lambda v) = \lambda (g v) \quad \text { for } v, v ^ {\prime} \text { in } T _ {x} (X), \lambda \in K.\tag{8}
\]

If \(x \in \mathbf{X}\) and \(u \in \mathrm{T}_g(\mathrm{G})\) , \(ux\) is by (4) the image of \(u\) under the tangent mapping at \(g\) to the mapping \(h \mapsto hx\) of \(\mathbf{G}\) into \(\mathbf{X}\) . Hence

\[
(9) \quad (u + u ^ {\prime}) x = u x + u ^ {\prime} x, \qquad (\lambda u) x = \lambda (u x) \quad \text { for } u, u ^ {\prime} \text { in } \mathrm{T} _ {g} (\mathrm{G}), \lambda \in \mathrm{K}.
\]

The above can be applied to the case of a Lie group operating on itself by left (resp. right) translation. The corresponding law of operation of T(G) on T(G) is defined by left (resp. right) translation of the Lie group T(G). Formulae (7), (8) and (9) are therefore valid in T(G).

PROPOSITION 8. Let \(\mathbf{G}_1\) and \(\mathbf{G}_2\) be Lie groups, \(\mathbf{X}_1\) and \(\mathbf{X}_2\) manifolds of class \(\mathbf{C}^r\) and \(f_i\) a law of left operation of class \(\mathbf{C}^r\) of \(\mathbf{G}_i\) on \(\mathbf{X}_i\) ( \(i = 1, 2\) ). Let \(\phi\) be a morphism of \(\mathbf{G}_1\) into \(\mathbf{G}_2\) and \(\psi\) a \(\phi\) -morphism of \(\mathbf{X}_1\) into \(\mathbf{X}_2\) . Then \(\mathrm{T}(\psi)\) is a \(\mathrm{T}(\phi)\) -morphism of \(\mathrm{T}(\mathbf{X}_1)\) into \(\mathrm{T}(\mathbf{X}_2)\) .

\(f_{2}\circ (\phi \times \psi) = \psi \circ f_{1},\) whence

\[
\mathrm{T} (f _ {2}) \circ (\mathrm{T} (\phi) \times \mathrm{T} (\psi)) = \mathrm{T} (\psi) \circ \mathrm{T} (f _ {1}).
\]

Let G be a Lie group, X a manifold of class \(C^{r}\) and \((g, x) \mapsto gx\) a law of left operation of class \(C^{r}\) of G on X. Let I be an open subset of K containing 0 and \(\gamma: I \to G\) a mapping of class \(C^{r}\) such that \(\gamma(0) = e\) . Let

\[
a = \mathrm{T} _ {0} (\gamma) 1 \in \mathrm{T} _ {e} (\mathrm{G}).
\]

Let \(x \in X\) . Using (4), \(ax\) is the image under the tangent mapping to \(\lambda \mapsto \gamma(\lambda)x\) of the tangent vector 1 to I at 0. Hence the vector field \(x \mapsto ax\) on \(X\) is the vector field defined by the mapping \((\lambda, x) \mapsto \gamma(\lambda)x\) in the sense of Differentiable and Analytic Manifolds, R, 8.4.5.
% semantic-fragments: legacy-input-seam
\relax{}
